\section{Related Work} \label{sec:related_work}

\subsection{Vision-Language-Action Models}

Vision-Language-Action (VLA) models have emerged as powerful robotic policies~\cite{brohan2022rt, kim2025openvla, li2024cogact, zheng2024tracevla, zheng2025universal, xiao2025ava}. Pioneering methods like OpenVLA~\cite{kim2025openvla} rely on discrete autoregressive tokenization, limiting inference frequency and struggling with continuous, multi-modal action distributions~\cite{zhong2025survey}. While subsequent architectures like $\pi_0$~\cite{black2410pi0} utilize conditional flow matching~\cite{lipman2022flow} to generate continuous actions, they remain predominantly 2D-centric and reactive~\cite{bu2025univla, black2410pi0, hu2025video, zheng2024tracevla}, lacking the robust 3D geometric understanding required for complex manipulation. Furthermore, although SpatialVLA~\cite{qu2025spatialvla} and BridgeVLA~\cite{li2025bridgevla} integrate explicit 3D information and auxiliary tasks, they still lack a predictive understanding of 3D scene dynamics. To address this gap, our GeoPredict introduces a geometry-aware predictive model that directly forecasts 3D scene evolution using 3D Gaussian Splatting~\cite{kerbl20233d}, equipping the VLA model with predictive geometric priors.



\subsection{Future Prediction for Robotic Manipulation}

Recent works explore future prediction to enhance robotic policies~\cite{cen2025worldvla, zhang2023storm, hafnermastering, micheli2022transformers, hansen2023td, hafner2022deep}. Some directly forecast future observations like RGB images~\cite{wu2024unleashing, hu2025video}, depths~\cite{zhangdreamvla}, or pointmaps~\cite{liu2025geometry}, while approaches like Seer~\cite{tianpredictive} condition inverse dynamics on these forecasted visual states. However, these methods face two key limitations. First, popular approaches like SuSiE~\cite{black2023zero} or UniPi~\cite{du2023learning} often rely on a single prediction step, which struggles to capture complex physical dynamics and introduces time-consuming denoising that reduces control frequency. Second, these 2D-centric methods~\cite{wu2024unleashing} struggle with multi-view consistency and lack the explicit 3D geometry crucial for accurate manipulation. To address this gap, we introduce GeoPredict, a geometry-aware predictive model that forecasts scene evolution over a full horizon of $H$ steps in explicit 3D space, learning a spatially coherent and predictive 3D representation to condition action generations.



\subsection{3D Scene Representation} 

3D scene representation is a fundamental problem in computer vision, traditionally relying on explicit structures like meshes~\cite{lorensen1998marching}, voxels~\cite{wu20153d} and point clouds~\cite{qi2017pointnet}. Recent research has shifted to neural representations that learn continuous models of scene geometry and appearance. Pioneered by Neural Radiance Fields (NeRF)~\cite{mildenhall2021nerf}, these implicit models often require subsequent acceleration~\cite{muller2022instant,xu2022point,hu2022efficientnerf} due to costly training and rendering. More recently, 3D Gaussian Splatting (3DGS)~\cite{kerbl20233d,yan2024street,wu20244d,ren2024scube} has emerged as a highly effective explicit alternative. 3DGS models scenes as a collection of 3D Gaussians with learned attributes, leveraging differentiable rasterization for high-fidelity and remarkably efficient rendering. While recent works like GWM~\cite{lu2025gwm} apply 3DGS to robotic world modeling, directly forecasting these massive Gaussian attributes remains computationally demanding. To circumvent this burden, we leverage 3DGS as a predictive scene geometry module solely during training, aiming to endow the VLA policy with intrinsic, foresight-driven spatial reasoning ability.
